\documentclass[11pt]{amsart}
\usepackage[margin=1.1in]{geometry}
\usepackage{amsmath,amssymb,amsthm}
\usepackage{mathtools}
\usepackage{xcolor}
\usepackage[colorlinks=true,linkcolor=blue!50!black,citecolor=blue!50!black,urlcolor=blue!50!black]{hyperref}
\usepackage[capitalise,nameinlink]{cleveref}

\newtheorem{theorem}{Theorem}[section]
\newtheorem{lemma}[theorem]{Lemma}
\newtheorem{proposition}[theorem]{Proposition}
\newtheorem{corollary}[theorem]{Corollary}
\theoremstyle{remark}
\newtheorem{remark}[theorem]{Remark}
\newtheorem{example}[theorem]{Example}

\crefname{lemma}{Lemma}{Lemmas}
\crefname{theorem}{Theorem}{Theorems}
\crefname{corollary}{Corollary}{Corollaries}
\crefname{proposition}{Proposition}{Propositions}
\crefname{remark}{Remark}{Remarks}

\newcommand{\E}{\mathbb{E}}
\renewcommand{\Pr}{\mathbb{P}}
\newcommand{\Inf}{\mathrm{Inf}}
\newcommand{\Var}{\mathrm{Var}}
\newcommand{\bits}{\{0,1\}}
\newcommand{\eps}{\varepsilon}
\newcommand{\JS}{\mathrm{JS}}

\title[Entropy--influence for monotone DNFs of bounded width]{The Fourier entropy--influence inequality for monotone DNF formulas of bounded width}
\author{Yitzchak Shmalo}
\address{Einstein Institute of Mathematics, The Hebrew University of Jerusalem, Jerusalem, Israel}
\email{yitzchak.shmalo@gmail.com}
\date{24 September 2026}

\begin{document}

\begin{abstract}
The Fourier entropy--influence conjecture of Friedgut and Kalai asserts that the spectral entropy of a Boolean
function is at most a universal constant times its total influence. It is open even for monotone functions. We prove
it with constant $w+8.1$ for every monotone DNF formula and every monotone CNF formula of width $w$, with no
restriction on the number of terms or on the number of times a variable occurs; for width two, that is for the
indicator that a random set of vertices of a graph contains an edge, the constant is $9.41$. More generally the
bound holds for every class of monotone functions that is closed under restrictions and whose nonconstant members
have a minterm or a maxterm of size at most $w$. The proof works with the uniform measure on the set of zeros of
the function, on which every coordinate is a fair coin unless it is blocked. This gives a bound on the entropy lost
when one coordinate is restricted, in terms of its influence relative to the variance. It also bounds the entropy by
the influences measured relative to the variance, and it reduces the conjecture for monotone functions to monotone
functions all of whose influences are small compared with the variance. Together with the influential-variable theorem of
O'Donnell, Saks, Schramm and Servedio it gives the constant $\log_2\log_2s+7.85$ for monotone functions computed by
decision trees with $s$ leaves, and $\log_2m+8.85$ for monotone DNF formulas with $m$ terms. The same classes satisfy the stronger inequality of
O'Donnell and Tan, also under product measures with biases bounded away from $0$ and $1$. So the conjecture also holds
for read-once formulas whose leaves are monotone DNFs or CNFs of bounded width and whose gates are AND, OR, XOR or
threshold gates of any fan-in, or bounded-width monotone DNFs whose inputs have biases bounded away from $0$ and $1$.
\end{abstract}

\maketitle

\section{Introduction}

For $f\colon\bits^n\to\bits$ let $F=1-2f$, $\chi_T(x)=(-1)^{\sum_{i\in T}x_i}$, $\widehat F(T)=\E[F\chi_T]$ with
$x$ uniform, and $P(T)=\widehat F(T)^2$. By Parseval $P$ is a probability distribution on the subsets of $[n]$; a
random set $S$ with law $P$ is the \emph{spectral sample} of $f$. The \emph{spectral entropy} and the \emph{total
influence} are
\[
H(f)=\sum_T P(T)\log_2\frac1{P(T)},\qquad I(f)=\sum_T|T|\,P(T)=\sum_{v}\Inf_v(f),
\]
where $\Inf_v(f)=\Pr[f(x^{v\leftarrow1})\neq f(x^{v\leftarrow0})]=\sum_{T\ni v}P(T)$. Friedgut and
Kalai~\cite{FK96} conjectured that $H(f)\le C\,I(f)$ for a universal constant $C$. The conjecture is known for
symmetric functions and read-once decision trees~\cite{OWZ11}, for read-once formulas~\cite{OT13}, for random
linear threshold functions~\cite{CKKLS18}, for functions of very small influence or very large entropy~\cite{Sha19},
for $\delta$-tribes and functions with separation properties~\cite{GMP26}, and for read-$k$ DNF formulas and
regular functions with constants $O(\log k)$ and $15.1$~\cite{Shm26}. The best general bounds lose a factor:
$H(f)=O(I(f)+\sum_AP(A)|A|\log|A|)$~\cite{KKLMS20}, and $H(f)\le\sum_vh(\Inf_v)=O(I(f)+\sum_v\Inf_v\log(1/\Inf_v))$,
with $h$ the binary entropy, which follows from the subadditivity of entropy~\cite{KMS12} and is also proved
in~\cite{Han25}.

The conjecture is open for monotone functions. There it would give total influence $\Omega(\log^2 n)$ for every
monotone graph property on $n$ vertices with constant variance~\cite{FK96}, where the best bound known is
$\Omega(\log^2n/(\log\log n)^2)$~\cite{KKLMS20}. The monotone case of the proofs above rests on an order of the
coordinates by importance (regular functions), on disjointness (read-once formulas) or on bounded overlap of terms
(read-$k$ DNFs). The results here use none of these.

For $f$ monotone write $q=\Pr[f=0]$ and $m=\min(q,1-q)$; then $\Var(F)=4q(1-q)$ lies between $2m$ and $4m$. A
\emph{minterm} of $f$ is a minimal set $M$ with $f(\mathbf 1_M)=1$, a \emph{maxterm} a minimal set $M$ with
$f(\mathbf 1_{[n]\setminus M})=0$. Let
\[
\Psi(r)=\log_2\!\bigl(e\,(r+2\sqrt r+2)\bigr)\quad(r\ge0),\qquad c_w=\frac2{\ln2}+\frac{2e^{-1/2}}{\ln 2}+\Psi(2^{w+1}-4).
\]
Numerically $2/\ln2+2e^{-1/2}/\ln2=4.6355$, $c_1=7.078$, $c_2=9.400$, $c_3=10.466$, $c_4=11.421$, and $c_w\le w+8.1$
for every $w\ge1$.

\begin{theorem}\label{thm:width}
Let $\mathcal C$ be a class of monotone Boolean functions that is closed under restrictions of variables to
constants and such that every nonconstant $f\in\mathcal C$ has a minterm or a maxterm with at most $w$ elements. Then
$H(f)\le c_w\,I(f)$ for every $f\in\mathcal C$.
\end{theorem}

\begin{corollary}\label{cor:dnf}
Every monotone DNF formula whose terms have at most $w$ variables, and every monotone CNF formula whose clauses have
at most $w$ variables, satisfies $H(f)\le (w+8.1)\,I(f)$. For $w=2$ the constant is $9.41$: for every graph $G$ on the
vertex set $[n]$ and every $U\subseteq[n]$, the function
\[
f(x)=\bigvee_{u\in U}x_u\ \vee\bigvee_{uw\in E(G)}(x_u\wedge x_w)
\]
satisfies $H(f)\le 9.41\,I(f)$, and so does the indicator that the set of ones of $x$ is a vertex cover of $G$.
\end{corollary}

The class of DNFs of width $w$ is closed under restrictions: setting $x_v=0$ deletes the terms containing $v$ and
setting $x_v=1$ removes $v$ from them. Every term is an implicant, so a nonconstant such $f$ has a minterm of size at
most $w$; the CNF case is the dual. For the graph formula, $1-f$ is the indicator that the set of ones is an
independent set of $G$ avoiding $U$. For this class the read-$k$ theorem of~\cite{Shm26}, with $k$ the maximum
degree of $G$, gives $(30+5.6\log_2k)\,I(f)$. Another member of the
class of \cref{thm:width}: for a fixed graph $K$ with $e$ edges, the property that a graph on $n$ vertices contains a
copy of $K$, as a function of the $\binom n2$ edge indicators, is a monotone DNF of width $e$, so its spectral entropy
is at most $(e+8.1)$ times its total influence, uniformly in $n$.

\Cref{thm:width} follows from a bound on the restriction step at a single coordinate. For a coordinate $v$ let
$f_0,f_1$ be the restrictions $x_v=0,1$ and
\[
\Delta_v(f)=H(f)-\tfrac12\bigl(H(f_0)+H(f_1)\bigr).
\]
Since $I(f)=\Inf_v(f)+\frac12(I(f_0)+I(f_1))$, a class of functions closed under restrictions in which every
nonconstant member has a coordinate with $\Delta_v\le K\Inf_v$ satisfies $H\le KI$, by induction on the number of
variables; this is the induction used by Gonz\'alez, MacManus and Pereyra~\cite{GMP26}.

\begin{theorem}\label{thm:step}
Let $f$ be monotone and let $v$ be a coordinate with $\Inf_v(f)>0$. Then
\[
\Delta_v(f)\le \Inf_v(f)\Bigl(\frac2{\ln 2}+\frac{2e^{-1/2}}{\ln2}+\Psi\Bigl(\frac{4m}{\Inf_v(f)}-2\Bigr)\Bigr).
\]
\end{theorem}

For an arbitrary Boolean function the step is at most $(2/\ln2+\log_2(e/\Inf_v))\Inf_v$ when $\Inf_v\le\frac12$;
summed along a telescoping it gives the bound of~\cite{KMS12} up to the constant. For monotone functions the
logarithm is of the influence relative to the variance. It cannot be removed altogether (\cref{rem:log}). Summing the
step along the telescoping of~\cite{Han25} gives the following.

\begin{corollary}\label{cor:han}
For every monotone $f$,
\[
H(f)\le\sum_v\Inf_v(f)\Bigl(4.6355+\Psi\Bigl(\frac{4m}{\Inf_v(f)}-2\Bigr)\Bigr)\le 7.08\,I(f)+\sum_v\Inf_v(f)\log_2\frac{4m}{\Inf_v(f)}.
\]
\end{corollary}

For balanced functions this is no better than the bound $H(f)\le\sum_vh(\Inf_v)$. For biased functions it is
stronger: for the
conjunction of $n$ variables every influence equals $2m$ and the right side is $8.08\,I$.

\begin{corollary}\label{cor:reduce}
Let $0<\tau\le2$ and $K>0$. Suppose every monotone $g$ with $\Inf_v(g)<\tau\,m(g)$ for all $v$ has a coordinate with
$\Delta_v(g)\le K\Inf_v(g)$. Then every monotone $f$ satisfies $H(f)\le\max(K,4.6355+\Psi(4/\tau-2))\,I(f)$.
\end{corollary}

So the conjecture for monotone functions reduces to monotone functions whose influences are all small compared with
their variance. Balanced tribes, majority and monotone graph properties with constant variance are of this kind;
\cref{sec:remarks} explains why the argument gives nothing more there.

In \cref{thm:width} the large influence comes from a short minterm. Decision trees supply another: by the theorem of
O'Donnell, Saks, Schramm and Servedio~\cite{OSSS05}, a function computed by a decision tree has a coordinate whose
influence is at least the variance divided by the expected depth of the tree. For a Boolean function $g$ let $\bar
d(g)$ be the least expected depth, on a uniformly random input, of a decision tree computing $g$, and let $L(f)$ be the
largest $\bar d(g)$ over the restrictions $g$ of $f$. A tree may be assumed to read no coordinate twice on a path;
then a leaf at depth $k$ is reached with probability $2^{-k}$, and the expected depth is the entropy of the leaf
reached. Restricting a tree does not add leaves, so $L(f)\le\log_2s$ if $f$ is computed by a tree with $s$ leaves. Also $L(f)$ is at most the
depth of any tree for $f$ and, averaging over the random choices, at most the least worst-case expected number of
queries of a randomized decision tree for $f$.

\begin{theorem}\label{thm:dt}
Every nonconstant monotone $f$ satisfies
\[
H(f)\le\bigl(4.6355+\Psi(2L(f)-2)\bigr)\,I(f)\le\bigl(\log_2L(f)+7.85\bigr)\,I(f).
\]
In particular, a monotone function computed by a decision tree with $s\ge2$ leaves satisfies
$H(f)\le(\log_2\log_2s+7.85)\,I(f)$.
\end{theorem}

For Boolean functions with $I(f)\ge1$ computed by a decision tree of expected depth $d$, Wan, Wright and
Wu~\cite{WWW14} proved $H(f)\le12d\,I(f)$. For monotone functions \cref{thm:dt} replaces the factor $d$ by its
logarithm, at the price of a maximum over restrictions. Write $s(f)$ for the sensitivity of $f$.

\begin{corollary}\label{cor:dtcases}
Let $f$ be a nonconstant monotone function.
\begin{enumerate}
\item[(a)] If $f$ is a DNF formula with $m$ terms or a CNF formula with $m$ clauses, then
$H(f)\le(\log_2m+8.85)\,I(f)$.
\item[(b)] If the minterms of $f$ have at most $w_1$ elements and its maxterms at most $w_0$ elements, then
$H(f)\le(\log_2(w_0w_1)+7.85)\,I(f)$.
\item[(c)] $H(f)\le(2\log_2s(f)+7.85)\,I(f)$.
\item[(d)] If $f$ has Fourier degree $k$, then $H(f)\le(4\log_2k+9.85)\,I(f)$.
\end{enumerate}
\end{corollary}

For DNF formulas with $m$ terms, monotone or not, the read-$k$ theorem of~\cite{Shm26} gives
$(30+5.6\log_2m)\,I(f)$. Part (b) is logarithmic in the widths, where \cref{thm:width} is linear in the smaller one.

O'Donnell and Tan~\cite{OT13} introduced a stronger inequality that is preserved under disjoint composition. With
$H^+(f)=\sum_{A\neq\emptyset}P(A)\log_2(1/P(A))$ and $\Var(f)=1-P(\emptyset)$ it reads
\begin{equation}\label{eq:feiplus}
H^+(f)\le C\,\bigl(I(f)-\Var(f)\bigr),
\end{equation}
and it implies $H(f)\le\max(C,1/\ln2)\,I(f)$~\cite[Proposition~2.1]{OT13}.

\begin{theorem}\label{thm:plus}
Every $f$ in a class $\mathcal C$ as in \cref{thm:width} satisfies \eqref{eq:feiplus} with
\[
C=C^+_w=\max\Bigl(8.80,\ 2\Bigl(1.0615+\frac2{\ln2}+\Psi(2^{w+1}-4)\Bigr)\Bigr).
\]
Numerically $C^+_1=12.78$, $C^+_2=17.43$, $C^+_3=19.56$, and $C^+_w\le2w+14.8$ for every $w\ge1$. Every nonconstant
monotone $f$ satisfies \eqref{eq:feiplus} with
$C=\max\bigl(8.80,\,2(1.0615+2/\ln2+\Psi(2L(f)-2))\bigr)\le2\log_2L(f)+14.33$.
\end{theorem}

By the composition theorem of~\cite{OT13}, functions satisfying \eqref{eq:feiplus} can be combined along a read-once
formula whose gates satisfy the biased form of \eqref{eq:feiplus} at the means of their inputs. Under every product
measure this holds for AND, OR and XOR gates of any fan-in with constant $2+4\log_2e<7.78$, for threshold gates
$\mathbf 1[y_1+\cdots+y_r\ge t]$ of any fan-in with constant $760$~\cite{Shm26}, and for gates of arity at most $a$
with a constant $2^{O(a)}$~\cite{OT13}.

\begin{corollary}\label{cor:compose}
Let $F$ be a read-once formula and $f=F(f_1,\dots,f_s)$, where $f_1,\dots,f_s$ are monotone DNF or CNF formulas of
width at most $w$ on pairwise disjoint sets of variables. Suppose every gate of $F$ is an AND, OR, XOR or threshold gate
of any fan-in, or a monotone DNF or CNF formula of width at most $w$ whose inputs are balanced, that is, equal to $1$
with probability $\frac12$. Then $H(f)\le\max(C^+_w,760)\,I(f)$, and $H(f)\le C^+_w\,I(f)$ if there are no threshold
gates.
\end{corollary}

For instance, the conjunction of any number of monotone $2$-DNFs on pairwise disjoint sets of variables is a monotone
DNF of unbounded width, and it satisfies $H\le17.43\,I$; so does $G(\mathrm{Maj}_3(x^1),\dots,\mathrm{Maj}_3(x^m))$ for
every monotone $2$-DNF $G$ and disjoint blocks $x^1,\dots,x^m$, since the majority of three unbiased bits is
unbiased. By the last statement of \cref{thm:plus}, the leaves $f_j$ in \cref{cor:compose} may also be monotone
functions with $L(f_j)\le L$, with $C^+_w$ replaced by $\max(C^+_w,2\log_2L+14.33)$.

The argument also works under product measures whose biases stay away from $0$ and $1$. For a product measure $\mu$ let
$P_\mu(A)$ be the squared coefficients of $F$ in the orthonormal basis of $L^2(\mu)$, $I_\mu$ and $\Var_\mu$ the
corresponding total influence and variance, and $\sigma_A^2=\prod_{i\in A}4\mu(x_i=1)\mu(x_i=0)$ (\cref{sec:product}).

\begin{theorem}\label{thm:product}
Let $0<\eta\le\frac12$, let $\mu$ be a product measure with $\mu(x_i=1)\in[\eta,1-\eta]$ for every $i$, and let
$\mathcal C$ be as in \cref{thm:width}. Then every $f\in\mathcal C$ satisfies
\[
\sum_{A\neq\emptyset}P_\mu(A)\log_2\frac{\sigma_A^2}{P_\mu(A)}\le\bigl(14.24+2w\log_2(1/\eta)\bigr)\bigl(I_\mu(f)-\Var_\mu(f)\bigr).
\]
\end{theorem}

So in \cref{cor:compose} a gate that is a monotone DNF or CNF formula of width at most $w$ may have inputs whose biases
lie in $[\eta,1-\eta]$, with $C^+_w$ replaced by $\max(C^+_w,14.24+2w\log_2(1/\eta))$.

\section{The restriction step}

\subsection{Erasure and Jensen--Shannon parts}
Fix $v$. For $T\subseteq[n]\setminus\{v\}$ the sets $T$ and $T\cup\{v\}$ carry the weights $P(T)$ and $P(T\cup v)$,
and the restrictions have $\widehat{F_b}(T)=\widehat F(T)\pm\widehat F(T\cup v)$. Let $\phi(t)=t\log_2(1/t)$ and
$\mathcal E(X,Y)=\phi(X)+\phi(Y)-\phi(X+Y)=(X+Y)\,h\bigl(Y/(X+Y)\bigr)$, where $h$ is the binary entropy. The law of
$S\setminus\{v\}$ is the average of the spectral distributions of $F_0$ and $F_1$, so
\[
\Delta_v=B_v+\JS_v,\qquad B_v=\sum_{T\not\ni v}\mathcal E\bigl(P(T),P(T\cup v)\bigr)=H(\mathbf 1[v\in S]\mid S\setminus\{v\}),
\]
and $\JS_v$ is the Jensen--Shannon divergence of the two spectral distributions.

\begin{lemma}\label{lem:js}
For every Boolean function, $\JS_v\le\frac{2}{\ln2}\sum_{T\not\ni v}\min\bigl(P(T),P(T\cup v)\bigr)\le\frac2{\ln2}\Inf_v$.
\end{lemma}

\begin{proof}
Put $u=\widehat F(T)$ and $w=\widehat F(T\cup v)$. The pair contributes
$J=\phi(u^2+w^2)-\frac12\phi((u+w)^2)-\frac12\phi((u-w)^2)$, which is $1/\ln2$ times the entropy
$\mathrm{Ent}(g^2)=\E g^2\ln g^2-\E g^2\ln\E g^2$ of the function $g$ on the two-point space with the uniform measure
and values $|u+w|$, $|u-w|$. By the logarithmic Sobolev inequality of the two-point space~\cite{Gro75},
$\mathrm{Ent}(g^2)\le\frac12(|u+w|-|u-w|)^2=2\min(u^2,w^2)$.
\end{proof}

\subsection{The down-set identity}
Let $f$ be monotone, $D=\{x:f(x)=0\}$, $q=|D|2^{-n}>0$, and let $\mu$ be the uniform measure on $D$, with
$\tilde\mu(T)=\E_\mu\chi_T$. Since $F=2\cdot\mathbf 1_D-1$,
\begin{equation}\label{eq:fmu}
\widehat F(T)=2q\,\tilde\mu(T)\qquad(T\neq\emptyset).
\end{equation}
Fix $v$. For $y\in\bits^{[n]\setminus\{v\}}$ call $y$ \emph{blocked} if $y\in D$ and $y+v\notin D$, and \emph{free}
if $y+v\in D$; here $y+v$ is $y$ with the $v$-th coordinate set to $1$. As $D$ is a down-set, the fibre of $D$ over a
blocked $y$ is $\{y\}$ and over a free $y$ is $\{y,y+v\}$. For $T\subseteq[n]\setminus\{v\}$ put
\[
a(T)=\frac1{|D|}\sum_{y\ \mathrm{blocked}}\chi_T(y),\qquad b(T)=\frac2{|D|}\sum_{y\ \mathrm{free}}\chi_T(y).
\]

\begin{lemma}\label{lem:downset}
For $T\subseteq[n]\setminus\{v\}$, $\tilde\mu(T\cup v)=a(T)$ and $\tilde\mu(T)=a(T)+b(T)$. With
$\eps=\Inf_v(f)/(4q)$,
\[
q\sum_T a(T)^2=\eps,\qquad q\sum_T b(T)^2=1-2\eps,\qquad 0\le\eps\le\tfrac12 .
\]
\end{lemma}

\begin{proof}
On a free fibre $x_v$ takes both values with equal weight, so $\chi_{T\cup v}$ averages to $0$ there, while on a
blocked fibre $x_v=0$; this gives both identities. Let $N_B,N_F$ be the numbers of blocked and free $y$. Then
$|D|=N_B+2N_F$, and $\Inf_v=N_B/2^{n-1}$ because the edges of the cube in direction $v$ on which $f$ changes are
exactly the pairs $\{y,y+v\}$ with $y$ blocked. Parseval on $\bits^{n-1}$ gives
$\sum_Ta(T)^2=2^{n-1}N_B/|D|^2$ and $\sum_Tb(T)^2=2^{n-1}\cdot4N_F/|D|^2$; multiplying by $q=|D|/2^n$ gives
$N_B/(2|D|)=\eps$ and $2N_F/|D|=1-2\eps$.
\end{proof}

Under $\mu$ a coordinate is a fair coin on free fibres and $0$ on blocked ones, and $2\eps$ is the $\mu$-probability
that it is blocked. By \eqref{eq:fmu}, $P(T\cup v)=4q^2a(T)^2$ for every $T$ and $P(T)=4q^2(a(T)+b(T))^2$ for $T$
nonempty. Also $\widehat F(\{v\})=\Inf_v$ for monotone $f$, so $P(\{v\})=\Inf_v^2$.

\subsection{The erasure bound}
The function $\Psi$ is increasing and concave, since $r+2\sqrt r+2$ is concave and positive. Let
$c_0=\max_{0<x\le1}x\log_2(e/x^2)=2e^{-1/2}/\ln2=1.75008$.

\begin{lemma}\label{lem:erasure}
For monotone $f$ and $\Inf_v(f)>0$, $B_v\le\Inf_v\bigl(c_0+\Psi(4m/\Inf_v-2)\bigr)$.
\end{lemma}

\begin{proof}
The dual $f^*(x)=1-f(\mathbf 1-x)$ is monotone, its coefficients are $-(-1)^{|T|}\widehat F(T)$, so it has the same
$P$, and $\Pr[f^*=0]=1-q$. Hence it suffices to prove the bound with $q$ in place of $m$; note $\Inf_v\le2q$. We use
$\mathcal E(X,Y)\le Y\log_2\bigl(e(X+Y)/Y\bigr)$, which follows from $X\log_2(1+Y/X)\le Y\log_2e$.

For $T=\emptyset$, $P(\{v\})=\Inf_v^2$ and $P(\emptyset)+P(\{v\})\le1$, so this pair contributes at most
$\Inf_v^2\log_2(e/\Inf_v^2)\le c_0\Inf_v$.

Let $T\neq\emptyset$. If $a(T)=0$ the pair contributes $0$. Otherwise let $s=b(T)/a(T)$. Since $\mathcal E$ is
homogeneous of degree one, the pair contributes $4q\cdot q\,a(T)^2\cdot\psi(s)$, where
$\psi(s)=\mathcal E((1+s)^2,1)\le\log_2\bigl(e((1+s)^2+1)\bigr)\le\Psi(s^2)$. The weights $w_T=q\,a(T)^2$ have total
$W\le\eps$, and $\sum_Tw_Ts_T^2\le q\sum_Tb(T)^2=1-2\eps$ by \cref{lem:downset}. By Jensen's inequality for the
concave increasing $\Psi$,
\[
\sum_T w_T\Psi(s_T^2)\le W\,\Psi\Bigl(\frac{1-2\eps}W\Bigr)\le \eps\,\Psi\Bigl(\frac{1-2\eps}{\eps}\Bigr),
\]
where the last step holds because $W\mapsto W\Psi(c/W)$ is nondecreasing: its derivative is $\Psi(z)-z\Psi'(z)$ at
$z=c/W$, which is at least $\Psi(0)>0$ by concavity. As $(1-2\eps)/\eps=4q/\Inf_v-2$ and $4q\eps=\Inf_v$, the pairs
with $T\neq\emptyset$ contribute at most $\Inf_v\Psi(4q/\Inf_v-2)$.
\end{proof}

\Cref{thm:step} is \cref{lem:erasure} together with \cref{lem:js}.

\section{Short certificates and decision trees}

\begin{lemma}\label{lem:cert}
Let $f$ be monotone with a minterm $M$ of size $k\ge1$. Then some $v\in M$ has
$\Inf_v(f)\ge 4q/(2^{k+1}-2)$. If $f$ has a maxterm of size $k$, some $v$ in it has
$\Inf_v(f)\ge4(1-q)/(2^{k+1}-2)$.
\end{lemma}

\begin{proof}
The second statement is the first for $f^*$. If $k=1$, then $M=\{v\}$ and $f=x_v\vee g$ with $g$ not depending on
$x_v$, so $\Inf_v=\Pr[g=0]=2q$. Let $k\ge2$, $M=\{v_1,\dots,v_k\}$, and let $\delta$ be the largest of the numbers
$\eps_{v_j}=\Inf_{v_j}/(4q)$. Under $\mu$ the conditional probability that $x_{v_j}=1$ given the other coordinates is
$\frac12$ on free fibres and $0$ on blocked ones, and the blocked fibres have probability $2\eps_{v_j}$. Hence
$p_j=\mu(x_{v_2}=\dots=x_{v_j}=1)$ satisfies $p_1=1$ and $p_{j+1}\ge\frac12(p_j-2\delta)$, so that
$p_k\ge2^{1-k}-\delta(2-2^{2-k})$. When $x_{v_2}=\dots=x_{v_k}=1$ the coordinate $v_1$ is blocked, since adding it
completes $M$; hence $2\delta\ge2\eps_{v_1}\ge p_k$, and therefore $\delta\ge1/(2^{k+1}-2)$.
\end{proof}

For $k=2$ the bound reads $\eps\ge\frac16$ at one end of every edge; it is attained by a single edge, where both ends
are occupied with probability $\frac13$ under the uniform measure on independent sets.

\begin{proof}[Proof of \cref{thm:width}]
Let $f\in\mathcal C$ be nonconstant. By \cref{lem:cert} it has a coordinate with $\Inf_v\ge4m/(2^{w+1}-2)$, so
$4m/\Inf_v-2\le2^{w+1}-4$ and \cref{thm:step} gives $\Delta_v\le c_w\Inf_v$, as $\Psi$ is increasing. The restrictions
are in $\mathcal C$ and have fewer variables, and constants have $H=I=0$; induction completes the proof. For the
estimate $c_w\le w+8.1$ use $\Psi(r)\le\log_2e+\log_2(r+2\sqrt r+2)$ with $r=2^{w+1}-4$, which gives
$c_w\le w+1+\log_2e+4.6355+\log_2(1+2^{(1-w)/2})$.
\end{proof}

\begin{remark}\label{rem:mixed}
The induction needs the short minterm only where no other bound is available. Let $\mathcal B\subseteq\mathcal C$ be
classes of monotone functions closed under restrictions, suppose $H\le C_{\mathcal B}I$ on $\mathcal B$, and suppose
every nonconstant $f\in\mathcal C\setminus\mathcal B$ has a minterm or a maxterm with at most $w$ elements. Then
$H\le\max(c_w,C_{\mathcal B})\,I$ on $\mathcal C$. For example, let $\mathcal B$ be the monotone regular functions,
which include the monotone threshold functions and satisfy $H\le15.1\,I$~\cite{Shm26}, and let $\mathcal C$ consist
of the functions $g\vee h$ with $g\in\mathcal B$ and $h$ a monotone DNF of width at most $w$, on the same variables or
not. A restriction of $g\vee h$ is $g_\rho\vee h_\rho$; if $h_\rho=0$ it lies in $\mathcal B$, and otherwise it is
constant or has a term of $h_\rho$, of size at most $w$, as an implicant. Hence, for instance,
$\mathrm{Maj}_n(x)\vee\bigvee_{uw\in E(G)}(x_u\wedge x_w)$ satisfies $H\le15.1\,I$ for every graph $G$ on $[n]$.
\end{remark}

\begin{proof}[Proof of \cref{thm:dt}]
Let $g$ be a nonconstant restriction of $f$, let $L=L(f)$, and let $T$ be a decision tree for $g$ whose expected
depth on a uniformly random input is at most $L$. If $\delta_i$ is the probability that $T$ reads the coordinate $i$,
then $\sum_i\delta_i$ is the expected depth, and the theorem of~\cite{OSSS05} gives
\[
\Var(G)\le\sum_i\delta_i\Inf_i(g)\le L\max_i\Inf_i(g),
\]
where $G=1-2g$. As $\Var(G)\ge2m(g)$, some coordinate has $4m(g)/\Inf_v(g)-2\le2L-2$, and \cref{thm:step} gives
$\Delta_v(g)\le(4.6355+\Psi(2L-2))\Inf_v(g)$. The restrictions of $f$ form a class closed under restrictions, so the
induction of the introduction applies. For the second bound, $\Psi(2L-2)=\log_2(2e)+\log_2L+\log_2(1+\sqrt{2L-2}/L)$
and $\sqrt{2L-2}/L\le1/\sqrt2$ for $L\ge1$, so $4.6355+\Psi(2L-2)\le\log_2L+4.6355+2.4427+0.7716$.
\end{proof}

\begin{proof}[Proof of \cref{cor:dtcases}]
(a) Evaluate the terms one after another, reading the unread variables of the current term until one of them is $0$
or all of them are $1$. For each term at most two new variables are read in expectation, so this tree has expected
depth at most $2m$. A restriction of $f$ is a DNF formula with at most $m$ terms, so $L(f)\le2m$, and
$\log_2(2m)+7.85=\log_2m+8.85$. CNF formulas are handled in the same way, reading variables until one is $1$.
(b) The least depth of a decision tree for a Boolean function $g$ is at most $C_0(g)C_1(g)$~\cite{BdW02}, where $C_b(g)$
is the largest, over the inputs $x$ with $g(x)=b$, of the least number of coordinates whose values at $x$ force the
value $b$. For monotone $g$, $C_1(g)$ and $C_0(g)$ are the largest sizes of a minterm and of a maxterm, and neither
increases under restriction. Hence $L(f)\le w_0w_1$.
(c) For monotone $g$ the certificate complexity $\max(C_0(g),C_1(g))$ equals the sensitivity $s(g)$~\cite{BdW02}, and
the sensitivity does not increase under restriction. Hence $L(f)\le s(f)^2$.
(d) For monotone $g$ the sensitivity equals the block sensitivity, which is at most $2\deg(g)^2$~\cite{BdW02}, and the
degree does not increase under restriction. Hence $L(f)\le4k^4$, and $\log_2(4k^4)+7.85=4\log_2k+9.85$.
\end{proof}

\section{Relative influences and the reduction}

\begin{proof}[Proof of \cref{cor:han}]
Han's telescoping~\cite[(3.4)]{Han25} writes $H(f)=\sum_{k=1}^n\E\,\Delta_k(f_{x})$, where $f_x$ is $f$ with the
coordinates after $k$ fixed to $x$, uniformly random, and $\Delta_k$ is its restriction step at $k$. Each $f_x$ is
monotone. By \cref{thm:step},
$\Delta_k(f_x)\le 4.6355\Inf_k(f_x)+G\bigl(\Inf_k(f_x),m(f_x)\bigr)$ with $G(x,m)=x\,\Psi(4m/x-2)$. The function $G$ is
the perspective of the concave $\Psi$ composed with the affine map $(x,m)\mapsto(x,4m-2x)$, hence jointly concave, and
it is increasing in $m$. Jensen's inequality with $\E\Inf_k(f_x)=\Inf_k(f)$ and $\E\,m(f_x)\le m(f)$ gives the first
bound. The second follows from $\Psi(r)\le\log_2\bigl(2e(r+2)\bigr)$.
\end{proof}

\begin{proof}[Proof of \cref{cor:reduce}]
Let $C=\max(K,4.6355+\Psi(4/\tau-2))$ and argue by induction on the number of variables. A nonconstant monotone $f$
either has a coordinate with $\Inf_v\ge\tau m$, where \cref{thm:step} gives $\Delta_v\le C\Inf_v$, or satisfies the
hypothesis; in both cases restricting $v$ and using the induction hypothesis on $f_0,f_1$ gives $H(f)\le C\,I(f)$.
\end{proof}

Running the same descent on a hypothetical counterexample shows where it would have to live: if $H(f)>C\,I(f)$ with
$C>7.08$, then, keeping at each step the restriction with the larger excess $H-CI$, one reaches a restriction $g$ of
$f$ with $H(g)-CI(g)\ge H(f)-CI(f)$ and $\Inf_v(g)<2^{9.08-C}m(g)$ for every coordinate $v$.

\section{The strengthened inequality}\label{sec:plus}

For a coordinate $v$ put $a_T=\widehat F(T)$ and $d_T=\widehat F(T\cup v)$ for $T\subseteq[n]\setminus\{v\}$, let
$D_vF=\frac12\bigl(F(x^{v\leftarrow0})-F(x^{v\leftarrow1})\bigr)$, and let
\[
\Delta^+_v=H^+(f)-\tfrac12\bigl(H^+(f_0)+H^+(f_1)\bigr),\qquad
J(a,d)=\phi(a^2+d^2)-\tfrac12\bigl(\phi((a+d)^2)+\phi((a-d)^2)\bigr)\ge0 .
\]

\begin{lemma}\label{lem:plusstep}
For every Boolean function and every coordinate $v$,
\begin{gather*}
\Delta^+_v=\phi(d_\emptyset^2)+\sum_{T\neq\emptyset}\bigl(\mathcal E(a_T^2,d_T^2)+J(a_T,d_T)\bigr),\\
\bigl(I-\Var\bigr)(f)=\Var(D_vF)+\tfrac12\bigl((I-\Var)(f_0)+(I-\Var)(f_1)\bigr).
\end{gather*}
\end{lemma}

\begin{proof}
The nonempty sets are the $T\neq\emptyset$ with $v\notin T$ and the sets $T\cup v$, including $\{v\}$; the restriction
$F_b$ has the coefficients $a_T\pm d_T$, and $\mathcal E(a^2,d^2)+J(a,d)=\phi(a^2)+\phi(d^2)-\frac12(\phi((a+d)^2)+
\phi((a-d)^2))$. For the second identity, $I(f)=\Inf_v+\frac12(I(f_0)+I(f_1))$, the means of $F_0,F_1$ are
$a_\emptyset\pm d_\emptyset$, so that $\Var(f)=\frac12(\Var(f_0)+\Var(f_1))+d_\emptyset^2$, and
$\Inf_v-d_\emptyset^2=\E(D_vF)^2-(\E D_vF)^2$.
\end{proof}

Consequently, a class closed under restrictions in which every nonconstant member has a coordinate with
$\Delta^+_v\le C\Var(D_vF)$ satisfies \eqref{eq:feiplus}, by induction on the number of variables.

\begin{lemma}\label{lem:neardict}
For monotone $f$ and every $v$ with $\Inf_v>0$,
$\Delta^+_v\le\Var(D_vF)\bigl(4/\ln2+\log_2(e(1+1/\Inf_v))\bigr)$. In particular $\Delta^+_v\le8.80\Var(D_vF)$
when $\Inf_v\ge\frac12$.
\end{lemma}

\begin{proof}
For monotone $f$, $D_vF$ takes the values $0$ and $1$, so $d_\emptyset=\Inf_v$ and
$\sum_{T\neq\emptyset}d_T^2=\Var(D_vF)=\Inf_v(1-\Inf_v)$. The function $A=\frac12(F_0+F_1)$ has
$\widehat A(T)=a_T$ and vanishes exactly where $v$ is pivotal, so $\sum_{T\neq\emptyset}a_T^2\le\E A^2=1-\Inf_v$. First,
$\phi(\Inf_v^2)=\frac2{\ln2}\Inf_v^2\ln\frac1{\Inf_v}\le\frac2{\ln2}\Var(D_vF)$. Next let
$M=\sum_{T\neq\emptyset}\min(a_T^2,d_T^2)\le\Var(D_vF)$ and $S=\sum_{T\neq\emptyset}(a_T^2+d_T^2)$. By the pointwise
bound in the proof of \cref{lem:js}, $\sum_{T\neq\emptyset}J(a_T,d_T)\le\frac2{\ln2}M$. Since
$\mathcal E(X,Y)\le\min(X,Y)\log_2\bigl(e(X+Y)/\min(X,Y)\bigr)$, the log-sum inequality gives
\[
\sum_{T\neq\emptyset}\mathcal E(a_T^2,d_T^2)\le M\log_2(eS/M),
\]
and as $x\log_2(eS/x)$ increases for $x\le S$,
this is at most $\Var(D_vF)\log_2(eS/\Var(D_vF))$, where $S/\Var(D_vF)\le1+(1-\Inf_v)/\Var(D_vF)=1+1/\Inf_v$. For
$\Inf_v\ge\frac12$ the constant is $4/\ln2+\log_2(3e)<8.80$.
\end{proof}

\begin{proof}[Proof of \cref{thm:plus}]
Let $f\in\mathcal C$ be nonconstant and let $v$ be the coordinate of \cref{lem:cert}, so that
$4m/\Inf_v-2\le2^{w+1}-4$; for the last statement, let $f$ be a nonconstant restriction of the given function and $v$ the
coordinate found in the proof of \cref{thm:dt}, with $4m/\Inf_v-2\le2L-2$, and replace $2^{w+1}-4$ by $2L-2$ below. If
$\Inf_v\ge\frac12$, \cref{lem:neardict} gives $\Delta^+_v\le8.80\Var(D_vF)$. Otherwise
\cref{lem:plusstep} gives $\Delta^+_v\le\phi(\Inf_v^2)+\sum_{T\neq\emptyset}\mathcal E(a_T^2,d_T^2)+\JS_v$. The middle
sum is the part of $B_v$ bounded in the proof of \cref{lem:erasure} by $\Inf_v\Psi(4m/\Inf_v-2)$, and
$\JS_v\le\frac2{\ln2}\Inf_v$ by \cref{lem:js}. As $\phi(x^2)\le\frac2{e\ln2}\,x<1.0615\,x$ and
$\Var(D_vF)=\Inf_v(1-\Inf_v)\ge\frac12\Inf_v$, we get $\Delta^+_v\le C^+_w\Var(D_vF)$, and \cref{lem:plusstep}
completes the induction. The estimate $C^+_w\le2w+14.8$ follows as in the proof of \cref{thm:width}, and the estimate
$2\log_2L+14.33$ as in the proof of \cref{thm:dt}, since $2(1.0615+2/\ln2+2.4427+0.7716)<14.33$.
\end{proof}

\begin{proof}[Proof of \cref{cor:compose}]
By~\cite[Theorem~1]{OT13}, if $g_1,\dots,g_k$ on disjoint sets of variables satisfy \eqref{eq:feiplus} with constant $C$
and $G$ satisfies the biased form of \eqref{eq:feiplus} with constant $C$ under the product measure with the means of
the $g_j$, then $G(g_1,\dots,g_k)$ satisfies \eqref{eq:feiplus} with constant $C$. The leaves satisfy it with $C^+_w$ by
\cref{thm:plus}, AND, OR, XOR and threshold gates with $2+4\log_2e<C^+_w$ or $760$~\cite{Shm26}, and a monotone DNF or
CNF gate of width at most $w$ with balanced inputs needs the inequality only under the uniform measure, where
\cref{thm:plus} gives $C^+_w$. Induction on the depth of $F$ and \cite[Proposition~2.1]{OT13} give the claim.
\end{proof}

\section{Product measures}\label{sec:product}

Let $\mu$ be a product measure on $\bits^n$ with $q_i=\mu(x_i=1)\in(0,1)$, put $\sigma_i^2=4q_i(1-q_i)$ and
$\sigma_A^2=\prod_{i\in A}\sigma_i^2$, and let $\chi^\mu_A=\prod_{i\in A}(x_i-q_i)/\sqrt{q_i(1-q_i)}$ be the orthonormal
basis of $L^2(\mu)$. For $F=1-2f=\sum_A\widetilde F(A)\chi^\mu_A$ put $P_\mu(A)=\widetilde F(A)^2$,
$I_\mu(f)=\sum_A|A|P_\mu(A)$, $\Var_\mu(f)=1-P_\mu(\emptyset)$, and
\[
L_\mu(f)=\sum_{A\neq\emptyset}P_\mu(A)\log_2\frac{\sigma_A^2}{P_\mu(A)},\qquad R_\mu(f)=I_\mu(f)-\Var_\mu(f),
\]
so that the biased form of \eqref{eq:feiplus} reads $L_\mu(f)\le C\,R_\mu(f)$. Fix a coordinate $v$ and write $q=q_v$,
$\sigma^2=\sigma_v^2$, $\tau=\sqrt{q/(1-q)}$, and $a_T=\widetilde F(T)$, $d_T=\widetilde F(T\cup v)$ for $T\not\ni v$. The
restrictions $f_0,f_1$ to $x_v=0,1$, with $\mu$ on the other coordinates, have the coefficients $a_T-\tau d_T$ and
$a_T+d_T/\tau$; they are weighted with $1-q$ and $q$, and the mixture of their spectra is $a_T^2+d_T^2$ because the
cross terms cancel. Put $\Inf^\mu_v=\sum_Td_T^2$ and
\[
J_q(a,d)=\phi(a^2+d^2)-(1-q)\,\phi\bigl((a-\tau d)^2\bigr)-q\,\phi\bigl((a+d/\tau)^2\bigr)\ge0 .
\]

\begin{lemma}\label{lem:biasedstep}
For every Boolean function and every coordinate $v$,
\begin{gather*}
L_\mu(f)-(1-q)L_\mu(f_0)-q\,L_\mu(f_1)=d_\emptyset^2\log_2\frac{\sigma^2}{d_\emptyset^2}
+\sum_{T\neq\emptyset}\bigl(\mathcal E(a_T^2,d_T^2)+J_q(a_T,d_T)\bigr)-\log_2\frac1{\sigma^2}\sum_{T\neq\emptyset}d_T^2,\\
R_\mu(f)-(1-q)R_\mu(f_0)-q\,R_\mu(f_1)=\sum_{T\neq\emptyset}d_T^2 .
\end{gather*}
\end{lemma}

\begin{proof}
The factors $\sigma_A^2$ with $A\not\ni v$ are the same in $f$ and in its restrictions, and each set $T\cup v$ carries the
extra factor $\sigma^2$; with the mixture $a_T^2+d_T^2$ this gives the first identity, as in \cref{lem:plusstep}. For the
second, $I_\mu$ drops by $\Inf^\mu_v$ and $\Var_\mu$ by $d_\emptyset^2$.
\end{proof}

The Jensen--Shannon part is bounded by the logarithmic Sobolev constant of the biased two-point space, and the factor
$\sigma^2$ pays for it.

\begin{lemma}\label{lem:biasedjs}
In this notation $J_q(a,d)\le\frac{c(q)}{\ln2}\,d^2$ for all real $a,d$, where $c(q)=\ln((1-q)/q)/(1-2q)$ and
$c(\frac12)=2$. The constant $c(q)$ is best possible, and $c(q)\le2+\ln(1/\sigma^2)$.
\end{lemma}

\begin{proof}
Put $u=a-\tau d$ and $w=a+d/\tau$. Then $(1-q)u^2+qw^2=a^2+d^2$, so $J_q(a,d)$ is $1/\ln2$ times $\mathrm{Ent}_\nu(g^2)$
for the function $g$ with values $|u|,|w|$ on the two-point space with measure $\nu=(1-q,q)$, and
$d^2=q(1-q)(w-u)^2\ge\Var_\nu(g)$. The logarithmic Sobolev inequality of this space with the optimal constant of
Diaconis and Saloff-Coste~\cite[Theorem~A.1]{DSC96} gives $\mathrm{Ent}_\nu(g^2)\le c(q)\Var_\nu(g)$. When $u$ and $w$
have the same sign, $d^2=\Var_\nu(g)$, so the optimality of $c(q)$ in that inequality shows that it cannot be lowered
here. For the last claim put $x=|1-2q|<1$. Then $c(q)=\frac1x\ln\frac{1+x}{1-x}=2\sum_{k\ge0}\frac{x^{2k}}{2k+1}$ and
$2+\ln\frac1{\sigma^2}=2+\ln\frac1{1-x^2}=2+\sum_{k\ge1}\frac{x^{2k}}k$, and $\frac2{2k+1}\le\frac1k$ for $k\ge1$.
\end{proof}

For monotone $f$ let $D=\{f=0\}$, let $\mu_D$ be $\mu$ conditioned on $D$, and let $\mu'$ be the marginal of $\mu$ on
the coordinates other than $v$. Call $y$ blocked or free as in \cref{lem:downset}; the blocked $y$ are the points at which
$v$ is pivotal. Let $\pi=\mu'(v\text{ pivotal})$ and $\theta=\mu_D(y\text{ blocked})$.

\begin{lemma}\label{lem:biasederasure}
For monotone $f$ and $\pi>0$, $d_\emptyset^2=\sigma^2\pi^2$, $\Inf^\mu_v=\sigma^2\pi$, and
\[
\sum_{T\neq\emptyset}\mathcal E(a_T^2,d_T^2)\le\Inf^\mu_v\log_2\Bigl(e\Bigl(1+\frac{2(1-q)}q+\frac{2(1-\theta)}{q\theta}\Bigr)\Bigr).
\]
\end{lemma}

\begin{proof}
Let $\mathrm B$ and $\mathrm{Fr}$ be the sets of blocked and free $y$. Under $\mu_D$ the coordinate $x_v$ is $0$ over a
blocked $y$ and has the law $(1-q,q)$ over a free one, so with $\alpha(T)=\E_{\mu_D}[\chi^\mu_T\mathbf 1_{\mathrm B}]$ and
$\gamma(T)=\E_{\mu_D}[\chi^\mu_T\mathbf 1_{\mathrm{Fr}}]$ we get $\E_{\mu_D}\chi^\mu_{T\cup v}=-\tau\alpha(T)$ and
$\E_{\mu_D}\chi^\mu_T=\alpha(T)+\gamma(T)$. As $F=2\cdot\mathbf 1_D-1$, $d_T=-2\mu(D)\tau\alpha(T)$ for every $T$ and
$a_T=2\mu(D)(\alpha(T)+\gamma(T))$ for $T\neq\emptyset$. The marginal of $\mu_D$ on $y$ has density $(1-q)/\mu(D)$ on
$\mathrm B$ and $1/\mu(D)$ on $\mathrm{Fr}$ with respect to $\mu'$, so Parseval for $\mu'$ gives
$\sum_T\alpha(T)^2=(1-q)^2\pi/\mu(D)^2$ and $\sum_T\gamma(T)^2=\mu'(\mathrm{Fr})/\mu(D)^2$. Hence
$\Inf^\mu_v=4q(1-q)\pi$, and $\alpha(\emptyset)=(1-q)\pi/\mu(D)$ gives $d_\emptyset^2=\sigma^2\pi^2$. For $T\neq\emptyset$
with $\alpha(T)\neq0$ put $s_T=\gamma(T)/\alpha(T)$. By $\mathcal E(X,Y)\le Y\log_2(e(X+Y)/Y)$ and $(1+s)^2\le2+2s^2$, the
pair contributes at most $w_Tg(s_T^2)$, where $w_T=4\mu(D)^2\tau^2\alpha(T)^2$ and
$g(z)=\log_2\bigl(e(1+2/\tau^2+2z/\tau^2)\bigr)$ is concave and increasing. The weights have total $W\le\Inf^\mu_v$ and
$\sum_Tw_Ts_T^2\le4\tau^2\mu'(\mathrm{Fr})$, so Jensen's inequality and the monotonicity of $W\mapsto Wg(c/W)$, as in the
proof of \cref{lem:erasure}, bound the sum by $\Inf^\mu_v\,g(4\tau^2\mu'(\mathrm{Fr})/\Inf^\mu_v)$. Finally
$\mu(D)=(1-q)\pi+\mu'(\mathrm{Fr})$ and $\theta=(1-q)\pi/\mu(D)$ give $8\mu'(\mathrm{Fr})/\Inf^\mu_v=2(1-\theta)/(q\theta)$.
\end{proof}

\begin{lemma}\label{lem:biasedcert}
Let $f$ be monotone with a minterm $M$ of size $k$, and suppose $q_i\ge\eta$ for $i\in M$, where $0<\eta<1$. Then some
$v\in M$ has $\theta\ge\eta^{k-1}(1-\eta)$.
\end{lemma}

\begin{proof}
For $k=1$ every $y$ with $(y,0)\in D$ is blocked, so $\theta=1$. Let $k\ge2$, $M=\{v_1,\dots,v_k\}$, and let $\delta$ be
the largest of the values of $\theta$ at $v_1,\dots,v_k$. Under $\mu_D$ the conditional probability that $x_{v_j}=1$
given the other coordinates is $q_{v_j}\ge\eta$ over a free point and $0$ over a blocked one, so
$p_j=\mu_D(x_{v_2}=\dots=x_{v_j}=1)$ satisfies $p_1=1$ and $p_{j+1}\ge\eta(p_j-\delta)$, whence
$p_k\ge\eta^{k-1}-\delta\eta/(1-\eta)$. When $x_{v_2}=\dots=x_{v_k}=1$ the coordinate $v_1$ is blocked, so $\delta\ge p_k$,
and therefore $\delta\ge\eta^{k-1}(1-\eta)$.
\end{proof}

\begin{proof}[Proof of \cref{thm:product}]
The dual $f^*(x)=1-f(\mathbf 1-x)$, under the measure with the biases $1-q_i$, has the same $P_\mu$, $\sigma_A^2$, $I_\mu$
and $\Var_\mu$, and its restrictions correspond to those of $f$; so we may assume that the nonconstant $f\in\mathcal C$
has a minterm of size $k\le w$, and we restrict the coordinate $v$ of \cref{lem:biasedcert}, for which
$\theta\ge\eta^{k-1}(1-\eta)\ge\eta^{k-1}/2$. Write $\rho=\sum_{T\neq\emptyset}d_T^2=\sigma^2\pi(1-\pi)$. By
\cref{lem:biasedjs}, $\sum_{T\neq\emptyset}J_q(a_T,d_T)-\log_2(1/\sigma^2)\rho\le\frac2{\ln2}\rho$, and
$d_\emptyset^2\log_2(\sigma^2/d_\emptyset^2)=2\Inf^\mu_v\,\pi\log_2(1/\pi)$ by \cref{lem:biasederasure}.

If $\pi\le\frac12$, then $\Inf^\mu_v\le2\rho$ and $2\pi\log_2(1/\pi)\le\frac2{e\ln2}<1.0615$. In \cref{lem:biasederasure},
$\frac{2(1-q)}q+\frac{2(1-\theta)}{q\theta}\le\frac4{\eta\theta}\le8\eta^{-k}$, so the erasure part is at most
$\Inf^\mu_v(\log_2(9e)+k\log_2(1/\eta))$. With \cref{lem:biasedstep} the step of $L_\mu$ is at most
$\rho\bigl(2\cdot1.0615+2\log_2(9e)+2/\ln2+2k\log_2(1/\eta)\bigr)<(14.24+2k\log_2(1/\eta))\,\rho$.

If $\pi>\frac12$, then $2\Inf^\mu_v\pi\log_2(1/\pi)\le\frac2{\ln2}\Inf^\mu_v(1-\pi)=\frac2{\ln2}\rho$. For the erasure part
argue as in \cref{lem:neardict}: $\sum_{T\neq\emptyset}\mathcal E(a_T^2,d_T^2)\le\rho\log_2(eS/\rho)$ with
$S=\sum_{T\neq\emptyset}(a_T^2+d_T^2)$. The function $(1-q)F_0+qF_1$ has the coefficients $a_T$, equals $1-2q$ where $v$ is
pivotal and $\pm1$ elsewhere, so its variance is at most $4(1-\pi)$; hence $S/\rho\le1+4/(\sigma^2\pi)\le9/\sigma^2$, and
$\sigma^2\ge4\eta(1-\eta)\ge2\eta$. The step of $L_\mu$ is then at most
$\rho\bigl(4/\ln2+\log_2(9e/2)+\log_2(1/\eta)\bigr)<(9.39+\log_2(1/\eta))\,\rho$.

In both cases \cref{lem:biasedstep} bounds the step of $L_\mu$ by $(14.24+2w\log_2(1/\eta))$ times the step of $R_\mu$,
the restrictions of $f$ lie in $\mathcal C$, and the measure on their coordinates has biases in $[\eta,1-\eta]$; induction
on the number of variables completes the proof.
\end{proof}

\section{Remarks}\label{sec:remarks}

\begin{remark}\label{rem:log}
The logarithm in \cref{thm:step} is needed. For odd $n=2m'+1$ and $y^*$ with $m'$ ones, the monotone function
$g_n(y,z)=\mathrm{Maj}_n(y)\vee(z\wedge[y=y^*])$ has $\Pr[g_n=0]=\frac12-2^{-n-1}$ and $\Inf_z=2^{-n}$, so
$\log_2(4m/\Inf_z)<n+1$. For $T\subseteq[n]$ the pair $(T,T\cup z)$ has weights $u_T^2$ and $v_T^2$ with
$v_T=\pm2^{-n}\chi_T(y^*)$ and $u_T=\widehat{M}(T)-2^{-n}\chi_T(y^*)$, where $M=1-2\,\mathrm{Maj}_n$. For $|T|=2j+1$,
$|\widehat M(T)|=2^{1-n}\binom{2m'}{m'}\binom{m'}{j}/\binom{2m'}{2j}$, and log-concavity gives
$\binom{2m'}{2j}\le(m'+1)\binom{m'}{j}^2$; hence $|\widehat M(T)|\ge2^{-n}2^{m'}/(\sqrt{m'}(m'+1))$, which for
$n\ge11$ is at least $2^{1-n}$, so that $|u_T|\ge\frac12|\widehat M(T)|$ and $4^nu_T^2\ge2^n/((n-1)(n+1)^2)$. Using
$\mathcal E(X,Y)\ge Y\log_2(1+X/Y)$ on the $2^{n-1}$ sets of odd size gives
$B_z\ge2^{n-1}4^{-n}\log_2\bigl(2^n/((n-1)(n+1)^2)\bigr)\ge\frac12(n-3\log_2(n+1))\Inf_z$. The exact ratios
$B_z/\Inf_z$ for $n=3,5,\dots,17$ are $3.52,4.24,4.99,5.78,6.62,7.48,8.37,9.27$.

Numerically the coefficient $1$ of the logarithm is attained as well. Majority loses a factor $2$ because
$\widehat M$ vanishes on the sets of even size, where the pairs cost nothing. Replace $\mathrm{Maj}_n$ by the threshold
function $[\,|y|\ge k\,]$, so that $M=1-2[\,|y|\ge k\,]$, and $y^*$ by a point with $k-1$ ones. Then
$2^n\widehat M(T)=2K_{k-1}(|T|-1;n-1)$ for $T\neq\emptyset$,
where $K_j(x;N)=\sum_i(-1)^i\binom xi\binom{N-x}{j-i}$ is the Krawtchouk polynomial, and the pairs can be summed exactly
in terms of $|T|$ and $|T\cap\operatorname{supp}y^*|$. For the best $k$ and $n=21,31,41,61,81$ the ratio $\Delta_z/\Inf_z$ is
$20.64,29.85,39.39,58.79,78.41$, against $\log_2(4m/\Inf_z)=20.62,30.17,41.09,60.29,80.58$; the bound of
\cref{thm:step} exceeds it by $6.1$ to $8.3$.
\end{remark}

\begin{proposition}\label{prop:maxinf}
For every $K$ there is a monotone function $f$ with a unique coordinate $z$ of largest influence and
$\Delta_z(f)>K\,\Inf_z(f)$. For instance, if every $y_j$ in the function $g_5$ of \cref{rem:log} is replaced by the majority
of $109$ new variables, the resulting monotone function has $546$ variables, $z$ is its unique coordinate of largest
influence, and $\Delta_z/\Inf_z=6.6596$.
\end{proposition}

\begin{proof}
Let $g$ be a Boolean function of $y_1,\dots,y_r$ and $z$, let $h_1,\dots,h_r$ be balanced Boolean functions on pairwise
disjoint sets of new variables, and let $f=g(h_1,\dots,h_r,z)$. With $G=1-2g$ and $H_j=1-2h_j$, the function $F=1-2f$
is obtained from the Fourier expansion of $G$ by substituting $H_j$ for the character of $y_j$, and
$\widehat{H_j}(\emptyset)=0$. Hence $\widehat F\bigl(Z\cup\bigcup_{j\in A}U_j\bigr)=\widehat G(A\cup Z)\prod_{j\in A}\widehat{H_j}(U_j)$
for $A\subseteq[r]$, $Z\subseteq\{z\}$ and nonempty sets $U_j$ of variables of $h_j$, and all other coefficients vanish:
the spectral sample of $f$ is that of $g$ with each $y_j$ in it replaced by an independent spectral sample of $h_j$.
Therefore $\Inf_z(f)=\Inf_z(g)$, a variable $u$ of $h_j$ has influence $\Inf_{y_j}(g)\Inf_u(h_j)$, and
$H(f)=H(g)+\sum_j\Inf_{y_j}(g)H(h_j)$. The restrictions of $f$ at $z$ are $g_b(h_1,\dots,h_r)$, so the same formula holds
for them, and since $\Inf_{y_j}(g)=\frac12(\Inf_{y_j}(g_0)+\Inf_{y_j}(g_1))$ the terms with $H(h_j)$ cancel in
$\Delta_z$: $\Delta_z(f)=\Delta_z(g)$.

Take $g=g_n$ from \cref{rem:log} and $h_j=\mathrm{Maj}_k$ with $k$ odd; then $f$ is monotone. As $k\to\infty$ the
influences of $\mathrm{Maj}_k$ tend to $0$, while $\Inf_z(f)=2^{-n}$, so for large $k$ the coordinate $z$ is the unique
coordinate of largest influence, and $\Delta_z(f)/\Inf_z(f)=\Delta_z(g_n)/\Inf_z(g_n)\ge B_z(g_n)/\Inf_z(g_n)
\ge\frac12(n-3\log_2(n+1))$ for $n\ge11$. For the instance, the influences of $g_5$ at the $y_j$ are at most $13/32$,
those of $\mathrm{Maj}_{109}$ equal $\binom{108}{54}2^{-108}<0.0766$, and $\frac{13}{32}\cdot0.0766<2^{-5}=\Inf_z$; the
ratio is that of $g_5$, computed exactly.
\end{proof}

With the threshold functions of \cref{rem:log} in place of $g_n$ the same construction gives monotone functions whose
unique coordinate $v$ of largest influence has $\Delta_v/\Inf_v$ within a few units of $\log_2(4m/\Inf_v)$ in the
computed range, so numerically the logarithm in \cref{thm:step} is attained with its coefficient even at such a
coordinate.

\begin{remark}
When every influence is small compared with $m$, $\eps$ is small at every coordinate, the free fibres carry almost
all of $\mu$, and \cref{lem:erasure} is no better than the bound $H(f)\le\sum_vh(\Inf_v)$. The argument used only the second moments of
$a$ and $b$. What can happen beyond them is shown by majority on an odd number of variables: $\widehat F$ vanishes on
sets of even size, so $a(T)=0$ when $|T|$ is odd and $a(T)+b(T)=0$ when $|T|$ is even and nonempty. Every pair then
has one weight equal to zero and $B_v=0$, although $\eps\sim n^{-1/2}$. A proof of the conjecture for monotone
functions along these lines would need an estimate of this alignment between the blocked and the free parts of the
down-set.

Numerically the step bound holds at every coordinate of largest influence: for every monotone function of at most six
variables, and for perfect matchings, tribes, majority and the property of containing a triangle, every coordinate $v$
of largest influence has $\Delta_v\le6.01\,\Inf_v$, the largest ratio being that of a perfect matching of eight
edges; over each of these classes the largest ratio at such coordinates equals the largest value of
$\min_v\Delta_v/\Inf_v$. For a perfect matching of $k$ edges the spectrum is a product over the edges, and the ratio,
computed exactly, increases to $6.0037$ at $k=8$ and then decreases to $5.735$ as $k\to\infty$. Among disjunctions of
$k\le40$ disjoint copies of any monotone function of at most four variables, which have the same product structure, the
largest ratio is again this one, and a search by simulated annealing over monotone functions of eight to twelve
variables found nothing above $5.72$. These searches do not reach the functions of \cref{prop:maxinf}, which need
hundreds of variables: in general the coordinate of largest influence does not work, and a proof has to choose the
coordinate by more than the size of its influence. The averaged form in \cref{prop:local} below needs no choice. On
every function computed $\sum_vB_v\le3.61\,I$; the largest value over the monotone functions of six variables is
$3.03$, and $3.61$ is again the perfect matching of eight edges.
\end{remark}

The erasure terms carry the whole conjecture, and in a form that involves one step of the spectral sample at a time.
Write $\log_2^+t=\max(\log_2t,0)$ and
\[
R(f)=\sum_{S:\,P(S)>0}P(S)\sum_{v\in S}\log_2^+\frac{P(S\setminus\{v\})}{P(S)} .
\]

\begin{proposition}\label{prop:local}
For every Boolean function $f$,
\[
R(f)\le\sum_vB_v\le\min\bigl(H(f),\,R(f)+\log_2(2e)\,I(f)\bigr).
\]
Consequently, if $\mathcal C$ is a class of Boolean functions closed under restrictions, then $H\le CI$ on $\mathcal C$
implies $R\le\sum_vB_v\le CI$ on $\mathcal C$, and $R\le KI$ on $\mathcal C$ implies
$H\le(K+\log_2(2e)+2/\ln2)\,I$ on $\mathcal C$.
\end{proposition}

\begin{proof}
By the chain rule, $H(f)=\sum_vH(\mathbf 1[v\in S]\mid\mathbf 1[u\in S],\,u<v)$, and each term is at least
$H(\mathbf 1[v\in S]\mid S\setminus\{v\})=B_v$. For $X\ge0$ and $Y>0$,
\[
Y\log_2^+(X/Y)\le Y\log_2\frac{X+Y}Y\le\mathcal E(X,Y)\le Y\log_2\frac{e(X+Y)}Y\le Y\bigl(\log_2(2e)+\log_2^+(X/Y)\bigr),
\]
and $\mathcal E(X,0)=0$. Applying this to $X=P(T)$, $Y=P(T\cup v)$ and summing over $v$ and $T\not\ni v$, that is over
the pairs $(S,v)$ with $v\in S=T\cup v$, gives the two bounds on $\sum_vB_v$. If $R\le KI$ on $\mathcal C$, then every
nonconstant $f\in\mathcal C$ has $\sum_vB_v\le(K+\log_2(2e))\sum_v\Inf_v$, and $B_v=0$ when $\Inf_v=0$, so some $v$
with $\Inf_v>0$ has $B_v\le(K+\log_2(2e))\Inf_v$; with \cref{lem:js} this gives
$\Delta_v\le(K+\log_2(2e)+2/\ln2)\Inf_v$, and the induction of the introduction applies.
\end{proof}

With $\mathcal C$ the class of all Boolean functions, the conjecture is thus equivalent to $R\le KI$: removing one
element of the spectral sample should increase its weight by at most a bounded factor, on average in the logarithm.
For monotone $f$ and $S=T\cup\{v\}$ with $T\neq\emptyset$ the ratio is $P(T)/P(S)=\bigl(1+b(T)/a(T)\bigr)^2$, and
\cref{lem:erasure} bounds its logarithmic average using only the second moments of $a$ and $b$.

\begin{remark}
The constant of \cref{thm:product} grows like $w\log_2(1/\eta)$, and the argument gives nothing as the biases tend to
$0$ or $1$: \cref{lem:biasederasure} uses only second moments, and under a product measure with small biases every
coordinate of a large matching has small relative influence. The restriction step is not what fails there. For the disjunction of $k$
disjoint conjunctions of two inputs, all of bias $q$, with $k\approx\ln2/q^2$ so that the output is balanced, the
biased spectral sample is a product over the pairs and can be computed exactly: at every input $v$ the conditional
entropy $H(\mathbf 1[v\in S]\mid S\setminus\{v\})$ is $o(\Inf^\mu_v)$ as $q\to0$, where
$\Inf^\mu_v=\sum_{A\ni v}P_\mu(A)$. The loss is in the second-moment estimate of \cref{lem:biasederasure}, which does not
see that each pair of weights $(P_\mu(T),P_\mu(T\cup v))$ is lopsided. Numerically the induction on coordinates itself
goes through: with $L_\mu$ and $R_\mu$ as in \cref{sec:product}, every
monotone $2$-DNF gate tested (fourteen graphs on up to six inputs, $261$ bias vectors each) has a coordinate at which
$L_\mu(G)-\E_bL_\mu(G_b)\le(2+4\log_2e)\bigl(R_\mu(G)-\E_bR_\mu(G_b)\bigr)$, the restriction $G_b$ being taken with probability $\mu(y_v=b)$;
the extreme case is a single edge with both biases near $1$, and on random $2$- and $3$-DNF gates with up to eight inputs,
with a local search on the biases, the largest ratio is $6.87$. By \cref{lem:biasedstep,lem:biasedjs}, what is missing
for a constant independent of the biases is a bound on the erasure part $\sum_T\mathcal E(a_T^2,d_T^2)$ that sees the
lopsided pairs. For a monotone $2$-DNF the uniform measure on the zeros, taken under a product
measure with biases $q_v$, is the hard-core model with activities $q_v/(1-q_v)$. Numerically the biased form holds
with the constant $2+4\log_2e$ of AND and OR. Over all monotone functions of at most five variables, and over random
monotone $2$- and $3$-DNFs with up to ten variables, with the biases chosen by local search, no ratio of the two sides
exceeds it; values above it computed in floating point at biases near $10^{-7}$ fall below it in $60$-digit
arithmetic.
\end{remark}

\begin{remark}\label{rem:numerics}
The statements were checked numerically. \Cref{lem:downset} holds to within $10^{-15}$, and
\cref{thm:step,cor:han} hold, for all $7{,}828{,}354$ monotone functions of six variables. On monotone DNF and CNF
formulas of widths two to four with up to $20$ variables and $120$ terms (paths, cycles, grids, hypercubes, complete,
complete bipartite, crown, Paley, random regular and multipartite graphs, and random formulas), the largest ratio of
$\Delta_v$ to the bound of \cref{thm:step} was $0.68$, and the lower bound of \cref{lem:cert} held in every case; a
perfect matching attains $\eps=\frac16$. The largest value of $H/I$ found for graph formulas was $4.52$, for a
perfect matching of ten edges. For the strengthened inequality, the identity of \cref{lem:plusstep} holds to within
$10^{-14}$ for all monotone functions of six variables, the two bounds used in the proof of \cref{thm:plus} are
attained to at most $0.75$ and $0.79$ of their values, and $H^+/(I-\Var)$ stays below $0.47\,C^+_w$ for DNFs of width
$w$. On the formula families above $H^+/(I-\Var)\le6.3$, and on conjunctions, disjunctions and majorities of two to
four monotone $2$-DNFs with up to $20$ variables $H/I\le4.27$. For \cref{thm:dt}, the least expected depth of a
decision tree was computed by dynamic programming over subcubes for all monotone functions of at most five variables and
for random monotone DNFs with eight to ten variables: the influence bound of~\cite{OSSS05} used in the proof held in
every case, with equality for $87$ of the $7{,}579$ nonconstant functions of five variables, and $H/I$ stayed below
$0.44$ times the bound of \cref{thm:dt}. On $400$ random monotone DNF formulas with up to $10$ variables and $29$
terms, $L(f)\le2m$ held with largest ratio $0.992$. The quantity $R$ of \cref{prop:local} is small in every case computed: over
all monotone functions of six variables the largest value of $R/I$ is $1.34$, against $3.74$ for $H/I$; majority has
$R=0$; and over the families above the largest value is $1.85$, for a perfect matching of eight edges. The scalar
inequalities in the proof of \cref{lem:erasure} (concavity and monotonicity of $\Psi$, monotonicity of
$W\mapsto W\Psi(c/W)$, and the resulting Jensen bound) and the arithmetic in the proof of \cref{lem:cert} were checked
in Lean~4 with Mathlib.
\end{remark}

\section*{Funding}
The research presented in this paper was supported by the European Research Council (ERC) under the European
Union's Horizon Europe research and innovation programme (grant agreement No.~101041711), by the Simons Foundation
as part of the Collaboration on the Mathematical and Scientific Foundations of Deep Learning, by Heights Labs, by
Convex Nexus Capital, by the Israel Science Foundation (grant number 2258/19), and by the Israel Science Foundation
(ISF Grant 4101/25).

\section*{Declaration of competing interest}
The author holds equity in Heights Labs and in Convex Nexus Capital and has a management role at Convex Nexus
Capital. These entities are included in the funding declaration above.

\section*{Declaration of generative AI and AI-assisted technologies}
The theorems, proofs and numerical computations in this paper were produced by the AI system Claude (Anthropic),
working within the author's math-brain project. The example in \cref{rem:log} was found earlier in the same project
by another group of Claude agents. The author takes full responsibility for the content of the paper.

\begin{thebibliography}{99}

\bibitem{BdW02} H.~Buhrman and R.~de~Wolf, Complexity measures and decision tree complexity: a survey,
\emph{Theoret. Comput. Sci.} 288 (2002), 21--43.

\bibitem{CKKLS18} S.~Chakraborty, S.~Karmalkar, S.~Kundu, S.~V.~Lokam and N.~Saurabh, Fourier entropy-influence
conjecture for random linear threshold functions, \emph{LATIN 2018}, Lecture Notes in Comput. Sci. 10807,
Springer, 2018, 275--289.

\bibitem{DSC96} P.~Diaconis and L.~Saloff-Coste, Logarithmic Sobolev inequalities for finite Markov chains,
\emph{Ann. Appl. Probab.} 6 (1996), 695--750.

\bibitem{FK96} E.~Friedgut and G.~Kalai, Every monotone graph property has a sharp threshold, \emph{Proc. Amer.
Math. Soc.} 124 (1996), 2993--3002.

\bibitem{GMP26} M.~J.~Gonz\'alez, P.~MacManus and M.~C.~Pereyra, Further evidence towards the Fourier
entropy-influence conjecture, arXiv:2606.00246 (2026).

\bibitem{Gro75} L.~Gross, Logarithmic Sobolev inequalities, \emph{Amer. J. Math.} 97 (1975), 1061--1083.

\bibitem{Han25} X.~Han, A new bound for the Fourier-entropy-influence conjecture, \emph{Combinatorica} 45 (2025).

\bibitem{KKLMS20} E.~Kelman, G.~Kindler, N.~Lifshitz, D.~Minzer and M.~Safra, Towards a proof of the
Fourier-entropy conjecture?, \emph{Geom. Funct. Anal.} 30 (2020), 1097--1138.

\bibitem{KMS12} N.~Keller, E.~Mossel and T.~Schlank, A note on the entropy/influence conjecture, \emph{Discrete
Math.} 312 (2012), 3364--3372.

\bibitem{OSSS05} R.~O'Donnell, M.~Saks, O.~Schramm and R.~A.~Servedio, Every decision tree has an influential
variable, \emph{46th Annual IEEE Symposium on Foundations of Computer Science (FOCS 2005)}, 31--39.

\bibitem{OT13} R.~O'Donnell and L.-Y.~Tan, A composition theorem for the Fourier entropy-influence conjecture,
\emph{ICALP 2013}, Lecture Notes in Comput. Sci. 7965, Springer, 2013, 780--791.

\bibitem{OWZ11} R.~O'Donnell, J.~Wright and Y.~Zhou, The Fourier entropy-influence conjecture for certain classes of
Boolean functions, \emph{ICALP 2011}, Lecture Notes in Comput. Sci. 6755, Springer, 2011, 330--341.

\bibitem{Sha19} G.~Shalev, On the Fourier entropy influence conjecture for extremal classes, arXiv:1806.03646 (2018).

\bibitem{Shm26} Y.~Shmalo, The Fourier entropy--influence inequality for read-$k$ DNF formulas and regular Boolean
functions, preprint (2026).

\bibitem{WWW14} A.~Wan, J.~Wright and C.~Wu, Decision trees, protocols and the entropy-influence conjecture,
\emph{Proceedings of the 5th Conference on Innovations in Theoretical Computer Science (ITCS 2014)}, 67--80.

\end{thebibliography}

\end{document}
